% ===================================================================
%  도표 제 10 호 — 바다. 항구.
%
%  Traduction, faite en 2026, en coreen d'aujourd'hui, sur l'ido de
%  texto/io. Regles de la colonne : voir 10-dopyo-01.tex. Une seule
%  scene, vingt-cinq alineas, cent renvois : les cles ne portent ni
%  c1 ni c2. Le titre du tableau tient dans le bloc « apar », comme
%  aux tableaux 1 a 6.
%
%  LE FAC-SIMILE IDO NE GRASSE NI « kaskado (10) » NI « kazino (21) »
%  — c'etait le tableau 9 ; ici il ne grasse pas « kaskado » non plus.
%  Les colonnes traduites suivent le francais et grassent, comme
%  toutes celles qui precedent.
%
%  LA COREE EST UNE PENINSULE, ET CE TABLEAU NE LUI COUTE RIEN. Trois
%  cotes, des ports de guerre a Jinhae et a Pyeongtaek, une marine
%  marchande parmi les premieres du monde : la langue a tout, et
%  souvent d'un mot de fond plutot que d'un compose savant —
%  잠수함 le sous-marin, 등대 le phare, 닻 l'ancre, 노 l'aviron,
%  돛 la voile, 그물 le filet, 밀물썰물 la maree, 반도 la presqu'ile,
%  해협 le detroit, 만 la baie, 곶 le cap, 부두 le quai, 방파제 la
%  digue, 갑판 le pont, 돛대 le mat, 키 le gouvernail, 뱃머리 la proue,
%  고물 la poupe.
%
%  Restent en anglais les objets que le coreen d'aujourd'hui nomme en
%  anglais, comme au tableau 10 telougou et pour la meme raison :
%  크루저, 프리깃, 어뢰, 보일러, 도크, 프로펠러, 크레인, 트럭, 전차,
%  카지노. Ce ne sont pas des trous — ce sont les mots que la langue
%  emploie.
%
%  TRENTE ET UNE PAIRES. Le tableau enumere des navires puis des
%  metiers de quai ; il attache moins qu'il ne compte, et le taux s'en
%  ressent — c'est le meme regime qu'au marche du tableau 15 telougou.
% ===================================================================

%%K t10-apar-1 apar
\VUsaut{4.0mm}
\VUtitre{15.0pt}{60}{도표 제 10 호}
\VUsaut{3.0mm}
\VUpk{11.4pt}{40}{바다. --- 항구.}
\VUsaut{3.0mm}

%%K t10-01-1 p
1. --- 여름에 어떤 이들은 산에서 고요와 서늘함을 찾지만, 다른 이들은 바닷가로 가기를 더
좋아한다. 우리도 지난해에 그렇게 했다. 우리는 \VUgras{큰 바다} \textsuperscript{(1)}를
아주 좋아하기 때문이다.

%%K t10-02-1 p
2. --- 나로 말하면, \VUgras{수평선} \textsuperscript{(2)}까지 펼쳐진 그 끝없는 물의 넓이를
바라보는 데서 큰 즐거움을 느낀다. 아메리카로 가는 손님들이 타는 거대한
\VUgras{기선} \textsuperscript{(3)}이 긴 \VUgras{바닷길}로 떠나는 것을 보는 데서도 그렇다.

%%K t10-03-1 p
3. --- 그래서 내 성미를 아시는 아버지는 방학을 보낼 곳으로 늘 우리 큰
\VUgras{군항} 가운데 하나에 가까운, 아주 조용한 모래밭을 고르신다.

%%K t10-03-2 p
지난번 머물 때 \VUgras{함대}가 거대한 \VUgras{방파제} \textsuperscript{(4)} 뒤에 닻을
내렸다. 그 방파제가 \VUgras{군항} \textsuperscript{(5)}을 막아 지켜 준다. 나는 바라던
대로 여러 \VUgras{군함}을 실컷 볼 수 있었다. 물속으로 잠겨 사라지는 작은
\VUgras{잠수함} \textsuperscript{(10)}, 그리고 거대한 \VUgras{장갑함} \textsuperscript{(9)}
까지. 그 배가 지금까지 두려워한 것은 거의 \VUgras{어뢰정} \textsuperscript{(8)}의 공격뿐이었다.

%%K t10-tit-2 sub
\VUsaut{3.0mm}
\VUpk{10.2pt}{40}{작은 배들.}
\VUsaut{3.0mm}

%%K t10-04-1 p
4. --- 그 배는 두꺼운 쇠 갑옷에서 약한 자리를 아주 \VUgras{솜씨 좋게} 찾아낼 줄 이미
알았다. 거기에 위험한 \VUgras{어뢰}를 박으려는 것이다. 그 폭발은 배를 잃게 만든다.
그래도 잠수함보다는 덜 무섭다. \VUgras{어뢰 막는 배}들이 그것을 쉽게 쫓기 때문이다.

%%K t10-05-1 p
5. --- 빠르고 갑옷을 두른 \VUgras{크루저} \textsuperscript{(6)}도 볼 수 있었다. 참된
장갑함만큼이나 다치지 않는 배다. 그리고 \VUgras{수송선} \textsuperscript{(12)} 몇 척,
\VUgras{포함} \textsuperscript{(7)} 서너 척, 끝으로 \VUgras{프리깃} \textsuperscript{(11)} 한
척이 있었다.
\VUgras{그 함대가 닻을 올리려고 움직일 때의 광경은 무엇보다 볼 만하다.}

%%K t10-06-1 p
6. --- 저 거대한 쇠 더미와, 맵시 있는 \VUgras{세 돛대 배}나 지금도 상선대에서 쓰이는
날렵한 \VUgras{돛단배} \textsuperscript{(13)} 사이의 짜임새는 얼마나 다른가! 내가
\VUgras{선창} \textsuperscript{(14)} 위를 거닐 때, 돛을 다 펴고 항구로 들어오는 그 배들을
보았다. 사실 \VUgras{바람}이 거스를 때면 그 배들은 제 이점을 많이 잃었다. 물웅덩이로
다시 들어오려면 돛을 다 내려야 했고, \VUgras{예인선} \textsuperscript{(16)} 한 척이 있어야
했다 — 그 배들을 데리러 가서 수수한 \VUgras{짐배} \textsuperscript{(17)}처럼 부두까지
끌어오는 배다.

%%K t10-tit-3 sub
\VUsaut{3.0mm}
\VUpk{10.2pt}{40}{상업 항구.}
\VUsaut{3.0mm}

%%K t10-07-1 p
7. --- 내가 말하는 이 항구에서 흥미로운 것은, 엄청난 공사로 손질한 군항만 있는 것이
아니라 큰 강 \VUgras{어귀}에 자리 잡은 훌륭한 \VUgras{상업 항구}도 함께 있다는 점이다.

%%K t10-08-1 p
8. --- 두 웅덩이를 가르는 땅 혀에는 방비가 되어 있다. 바다로 내민 \VUgras{포대}와
\VUgras{보루}가 그것을 지킨다. \VUgras{성벽} \textsuperscript{(15)} 위를 거닐려면 특별한
허락이 있어야 한다. 나는 그것을 삼촌 덕분에 쉽게 얻었다 — 삼촌은 \VUgras{조선소}
\textsuperscript{(18)}의 기술자이시다 — 그래서 넓은 지붕 아래에서 짓고 있는 배들을 보러
갔다.

%%K t10-09-1 p
9. --- 장갑함을 물에 내리는 것까지 보았다. 그 거대한 덩어리가 제 무게에 맡겨져 요람 같은
틀 위를 미끄러지기 시작하여 강물 위에 뜨던 그 순간, 온 무리가 느낀 감동을 나는 기억한다.

%%K t10-10-1 p
10. --- 조선소로 가려면 \VUgras{연병장} \textsuperscript{(19)}을 가로질러야 한다 — 거기에
주둔군 병사들이 날마다 훈련하러 온다 —, 그 다음에 쇠로 된 \VUgras{작은 다리}
\textsuperscript{(20)}를 건너야 한다. 그 다리가 연병장을 \VUgras{병기창} \textsuperscript{(21)}
과 이어 준다.

%%K t10-tit-4 sub
\VUsaut{3.0mm}
\VUpk{10.2pt}{40}{병기창과 도크.}
\VUsaut{3.0mm}

%%K t10-11-1 p
11. --- 삼촌과 함께 나는 그 큰 건물을 둘러보았다. \VUgras{대포} \textsuperscript{(22)}와
\VUgras{포탄} \textsuperscript{(23)}과 \VUgras{작렬탄}이 가득한 곳이다.
\VUgras{쇠 공장} \textsuperscript{(24)}도 보았다 — 거기에서 기선의 \VUgras{보일러}
\textsuperscript{(25)}를 만든다.

%%K t10-12-1 p
12. --- 아주 가까이에 \VUgras{도크} \textsuperscript{(26)} — 고치는 도크 — 가 있고, 아주
볼 만한 \VUgras{뜬 도크} \textsuperscript{(27)}도 있다. 거기에서 나는 고치러 온 배 위의
\VUgras{프로펠러} \textsuperscript{(28)}와 \VUgras{용골} \textsuperscript{(29)}과 \VUgras{키}
\textsuperscript{(30)}를 아주 가까이에서 볼 수 있었다.

%%K t10-13-1 p
13. --- 상업 항구도 군항 못지않게 볼 만하다. 나는 부두에서 여러 시간을 입을 벌린 채
서 있었다. 거기에는 강을 거슬러 오르는 \VUgras{외륜선} \textsuperscript{(31)}이 매여
있고, \VUgras{돛대 세우는 기중기} \textsuperscript{(32)}가 일하는 것도 보인다 — 그것은
거대한 짐을 깃털처럼 들어 올린다.

%%K t10-tit-5 sub
\VUsaut{4.0mm}
\VUpk{10.2pt}{40}{도선사.}
\VUsaut{3.0mm}

%%K t10-14-1 p
14. --- 정박지를 빠져나가는 \VUgras{대양 여객선} \textsuperscript{(34)}을 보고 감탄했다.
그 \VUgras{깃발} \textsuperscript{(35)}이 바람에 활짝 펄럭이고, 솜씨 있는 \VUgras{도선사}
\textsuperscript{(36)}가 배를 이끈다. 그는 \VUgras{부표} \textsuperscript{(40)}와
\VUgras{길잡이 말뚝}으로 표시한 좁은 \VUgras{물길}을 따라가는 법을 놀랍도록 잘 안다 —
\VUgras{거룻배} \textsuperscript{(41)}들을 피하면서. 그 배들은 네모난 돛 하나로 힘겹게
방향을 잡는다. 나는 \VUgras{앞부분} \textsuperscript{(37)} — \VUgras{뱃머리} — 에서
이등 \VUgras{선실} \textsuperscript{(39)}을, \VUgras{뒷부분} \textsuperscript{(38)} —
\VUgras{고물} — 에서 일등 손님용 큰 방을 알아보았다.

%%K t10-15-1 p
15. --- 이따금 \VUgras{짐배} \textsuperscript{(42)}에 짐 싣는 것도 보았다. 거기에는
\VUgras{창고} \textsuperscript{(43)}와 \VUgras{바다 정거장} \textsuperscript{(44)}의 물건을
방금 쌓아 놓았다. \VUgras{부두 철길} \textsuperscript{(45)}로 오는 숱한 기차가 그것을
날라 온 것이다.

%%K t10-tit-6 sub
\VUsaut{3.0mm}
\VUpk{10.2pt}{40}{노 젓는 즐거움.}
\VUsaut{3.0mm}

%%K t10-16-1 p
16. --- 나는 \VUgras{물웅덩이 도크} \textsuperscript{(53)}에서 자주 배를 저었다 — 거기로
\VUgras{작은 배} \textsuperscript{(46)}들이 피해 들어온다 —, \VUgras{노} \textsuperscript{(47)}
를 다루는 것이 나에게는 기쁨이었다. \VUgras{갑판} \textsuperscript{(48)} 위로 올라갔는데,
\VUgras{돛단배} \textsuperscript{(49)}의 갑판이다. \VUgras{굵은 밧줄} \textsuperscript{(52)}
을 잡고 \VUgras{돛대} \textsuperscript{(51)}를 타고 \VUgras{활대} \textsuperscript{(50)}까지
올라갔다. 그러고는 \VUgras{작은 배} \textsuperscript{(54)}로 나들이를 이어 갔다 — 묵직한
\VUgras{고깃배} \textsuperscript{(55)}들 사이로, 그 배들 위에서 \VUgras{그물}
\textsuperscript{(56)}이 마른다.

%%K t10-17-1 p
17. --- \VUgras{부두} \textsuperscript{(58)} 위에서 내 눈앞으로 온갖 \VUgras{짐}
\textsuperscript{(59)}이 줄지어 지나갔다 — \VUgras{궤짝} \textsuperscript{(60)}에 담기거나
\VUgras{꾸러미} \textsuperscript{(61)}로 묶인 것이다. 그것들이 큰 거룻배의 \VUgras{실은 짐}
\textsuperscript{(63)}이 되었다. 힘센 \VUgras{크레인} \textsuperscript{(62)}이 그것을 내려
\VUgras{트럭} \textsuperscript{(64)} 위에 실었고, 그동안 \VUgras{짐 나르는 사람}들이
\VUgras{세관} \textsuperscript{(65)}에서 그것을 신고하고 세금을 냈다.

%%K t10-18-1 p
18. --- 끝으로 \VUgras{도는 다리} \textsuperscript{(66)}나 \VUgras{갑문} \textsuperscript{(69)}
에 이르렀을 때 — \VUgras{준설선} \textsuperscript{(68)} 앞을 지나면서, 그 배가
\VUgras{운하} \textsuperscript{(67)}를 치운다 — 나는 \VUgras{신호탑} \textsuperscript{(70)}
가까이 선창에 내렸다.

%%K t10-19-1 p
19. --- 그 선창 반대쪽으로 \VUgras{함재정} \textsuperscript{(71)}들이 와서 댄다. 함대
장교들을 뭍으로 나르는 배다. 뱃사람들이 박자에 맞추어 노를 드는 것을 보는 일은 볼 만한
광경이다. 그동안 \VUgras{물결} \textsuperscript{(72)}에 실린 배는 내리는 곳 층계에 쉽게
닿는다. 이 자리에서는 \VUgras{모래밭} \textsuperscript{(73)}의 \VUgras{밀물썰물}과
\VUgras{밀물}, \VUgras{썰물}의 움직임을 더 잘 살필 수 있다.

%%K t10-tit-7 sub
\VUsaut{3.0mm}
\VUpk{10.2pt}{40}{장교 구락부.}
\VUsaut{3.0mm}

%%K t10-20-1 p
20. --- 집으로 돌아오려고 나는 \VUgras{전차} \textsuperscript{(74)}를 탔다. 그
\VUgras{정류장} \textsuperscript{(75)}은 장교 구락부 앞에 세운 \VUgras{기둥} 가까이에 있다.

%%K t10-20-2 p
구락부 \VUgras{발코니} \textsuperscript{(76)}에서 — \VUgras{닻} \textsuperscript{(77)}과
대포와 포탄으로 꾸민 곳이다 — 나는 해군 장교들의 훌륭한 모임을 자주 보았다.
\VUgras{중위} \textsuperscript{(78)}들은 제 \VUgras{긴 칼}을 차고, \VUgras{포병 대위}
\textsuperscript{(79)}와 \VUgras{보병 대위} \textsuperscript{(80)}는 제 \VUgras{군도}를
차고 있었다. 그 뒤에는 \VUgras{수병} \textsuperscript{(81)} 하나가 있어, 멀리 일어나는 것을
알아보고 싶을 때 \VUgras{망원경}을 건넸다.

%%K t10-21-1 p
21. --- 젊은 \VUgras{소위} \textsuperscript{(82)}들도 보았다. 그들은 제 \VUgras{함장}
\textsuperscript{(83)}이나 \VUgras{제독} \textsuperscript{(84)}에게서 지시를 받고 있었고,
그동안 \VUgras{병조장} \textsuperscript{(85)}이 그것을 배로 나르려고 기다리고 있었다.

%%K t10-22-1 p
22. --- 그 발코니에서 보이는 경치는 훌륭하다. \VUgras{천막} \textsuperscript{(86)}과
\VUgras{탈의실} \textsuperscript{(87)}이 있는 모래밭 전체가 한눈에 들어오고, 미역 감는
시간이면 \VUgras{미역 감는 사람} \textsuperscript{(88)}들이 \VUgras{카지노}
\textsuperscript{(89)} 앞에서 즐겁게 뛰노는 것이 보인다.

%%K t10-23-1 p
23. --- 모래밭 끝에서 바닷가는 바위투성이의 작은 \VUgras{반도} \textsuperscript{(91)}를
이룬다. 거기로 \VUgras{큰 물결} \textsuperscript{(92)}이 와서 부서지고, 그 위에
\VUgras{등대} \textsuperscript{(93)}를 세웠다. 밤이면 그것이 뱃사람들에게 길을 알려 준다.
그 반도가 뭍과 이어지는 것은 아주 좁은 \VUgras{잘록한 땅} \textsuperscript{(90)}을
통해서뿐이다.

%%K t10-tit-8 sub
\VUsaut{3.0mm}
\VUpk{10.2pt}{40}{바닷가의 두루 보기.}
\VUsaut{3.0mm}

%%K t10-24-1 p
24. --- \VUgras{벼랑} \textsuperscript{(94)}이 길게 뻗어 있는데, 바다가 그것을 갈라놓았다.
바다는 거기에 제멋대로 \VUgras{만} \textsuperscript{(95)}과 \VUgras{곶}을 잇달아 그려
놓아, 더할 나위 없이 그림 같게 만들었다.

%%K t10-24-2 p
\VUgras{대지} \textsuperscript{(96)} 위에는 — 그것이 \VUgras{해협} \textsuperscript{(98)}을
내려다본다 — 아주 중요한 \VUgras{요새} \textsuperscript{(97)} 하나와 「별장」 몇 채가 있다.

%%K t10-25-1 p
25. --- 그 해협이 뭍에서 아주 위험한 \VUgras{섬} \textsuperscript{(99)}을 갈라놓는다. 그
섬 둘레에는 \VUgras{암초} \textsuperscript{(100)}와 \VUgras{부서지는 물결}이 있어, 작은
배는 물론 큰 배까지도 이따금 얹혀 버린다. 구조가 빠르고 \VUgras{구명정}이 서둘러 와도
그 \VUgras{난파}는 무섭다. 춘분 가을분의 \VUgras{큰바람} 때에는 여러 배가 사람과 짐과
함께 사라졌다.

%%K t10-25-2 p
그런 이웃이 있어도 \VUgras{그 항구}에는 뱃사람들이 아주 많이 드나든다.

\VUsaut{6.0mm}
\VUfilet{22mm}
